\scanpage{30}
% archive sheet number: 29
\mspage{1}
\uline{Structure du groupe} $\mathrm{Sl}(2,\Lambda)$.

1. Soit $M$ un module libre de rang 2 sur un anneau $\Lambda$, avec une structure symplectique
\[
e_M\in(\bigwedge\nolimits^2M)^*.\tag*{(1)}
\]
La donnée d'un tore maximal déployé $T$ de $\mathrm{Sl}(M)$ \uline{revient} à la donnée d'une décomposition de $M$ en somme de deux sous-modules libres \inserted{$L'_i$} de rang 1, formant un ensemble \deleted{$\{L,L'\}$} à deux éléments $\varepsilon_T$ --- \unsure{Il|ce} \deleted{\unsure{leurs|leur}} sont les deux sous-modules propres de $M$, et on a
% slanted
\margin{Si $\operatorname{Spec}\Lambda$ connexe --- Sinon ce n'est le cas que \unsure{lorsqu'on|lorsqu'on a} \uline{\unsure{donne|donné}} un déploiement de $T$, i.e. une paire d'iso $T\simeq G_m$ opposés l'un de l'autre}
\[
\bigotimes_{i\in\varepsilon_T}L'_i\simeq(\bigwedge\nolimits^2M)\wedge_{\pm1}\varepsilon_T\simeq \Lambda\wedge_{\pm1}\varepsilon_T\tag*{(2)}
\]
en utilisant l'isomorphisme
\[
\bigwedge\nolimits^2 M\xrightarrow{\sim}\Lambda\qquad e_M\longmapsto 1\tag*{(3)}
\]
défini par la base $e_M$. L'opération de $T$ sur chaque $L_i$ définit un caractère
\[
T\xrightarrow{\chi_i}\mathbb G_m\qquad i\in\varepsilon_T\tag*{(4)}
\]
qui est un isomorphisme, et (2) implique que l'on a
\[
\prod_{i\in\varepsilon_T}\chi_i=1\tag*{(5)}
\]
i.e. les deux caractères $\chi_i$ sont inverses l'un de l'autre : ce sont les deux isomorphismes de $T$ sur $G_m$ (quand $\operatorname{Spec}\unsure{A|\Lambda}$ est connexe --- sinon on peut dire \inserted{encore} que ce sont les deux iso. ``canoniques'' de $T$ sur $G_m$). Choisir l'un des ces isomorphismes revient donc à choisir
% slanted
\margin{ils sont donnés avec le déploiement}
